\documentclass[11pt]{article}
\usepackage[a4paper,margin=1in]{geometry}
\usepackage{amsmath,amssymb,amsthm}
\usepackage{mathpartir}
\usepackage{listings}
\usepackage{hyperref}

\lstset{
  basicstyle=\ttfamily\small,
  columns=fullflexible,
  breaklines=true,
  keepspaces=true,
  showstringspaces=false,
}

\theoremstyle{plain}
\newtheorem{theorem}{Theorem}[section]
\newtheorem{lemma}[theorem]{Lemma}
\newtheorem{proposition}[theorem]{Proposition}
\theoremstyle{definition}
\newtheorem{definition}[theorem]{Definition}
\theoremstyle{remark}
\newtheorem{remark}[theorem]{Remark}

% ------------------------------------------------------------------
% Notation (Shoenfield, Mathematical Logic, Ch.6, with code(u) for
% the encoding in place of corner-quotes).
% ------------------------------------------------------------------
\newcommand{\code}[1]{\mathit{code}(#1)}
\newcommand{\sub}{\mathbf{sub}}
\newcommand{\sbf}{\mathbf{sbf}}
\newcommand{\sbt}{\mathbf{sbt}}
\newcommand{\sbftwo}{\mathbf{sbf_2}}
\newcommand{\sbttwo}{\mathbf{sbt_2}}
\newcommand{\subtwo}{\mathbf{sub_2}}
\newcommand{\thmT}{\mathbf{thmT}}
\newcommand{\num}{\mathbf{num}}
\newcommand{\Pair}{\mathbf{Pair}}
\newcommand{\seq}[1]{\langle #1 \rangle}
\newcommand{\Fun}{\mathit{Fun}}
\newcommand{\Sub}{\mathit{Sub}}
\newcommand{\Subtl}{\mathit{Subtl}}
\newcommand{\Term}{\mathit{Term}}
\newcommand{\AFor}{\mathit{AFor}}
\newcommand{\For}{\mathit{For}}
\newcommand{\Vble}{\mathit{Vble}}
\newcommand{\Fr}{\mathit{Fr}}
\newcommand{\Prf}{\mathit{Prf}}
\newcommand{\Prr}{\mathit{Pr}}
\newcommand{\Thmm}{\mathit{Thm}}
\newcommand{\Num}{\mathit{Num}}
\newcommand{\Con}{\mathit{Con}}
\newcommand{\Closed}{\mathit{Closed}}
\newcommand{\substT}{\mathit{substT}}
\newcommand{\substF}{\mathit{substF}}
\newcommand{\Len}{\mathit{Len}}
\newcommand{\SN}{\mathit{SN}}
\newcommand{\lh}{\mathit{lh}}

\title{A formalisation of G\"odel's Second Incompleteness Theorem\\for the Basic Recursive Arithmetic of Church and Guard}
\author{}
\date{}

\begin{document}
\maketitle

\begin{abstract}
We report a complete, axiom-free, machine-checked formalisation in Agda
of G\"odel's second incompleteness theorem for Church's \emph{basic
recursive arithmetic} (BRA), following the proof outline given by
R.~Guard in his 1963 paper.  Beyond the formalisation itself, the work
exposes a number of implicit assumptions, notational ambiguities, and
minor typographical issues in Guard's text that are worth recording
explicitly.  We also record an architectural simplification we have since
carried out: replacing the dedicated multi-variable simultaneous-substitution
machinery by nested single-variable substitution together with a
numeral-inertness lemma (the internalisation of ``a numeral is closed''),
which removes roughly ten thousand lines.  Throughout we follow Shoenfield's notation
(\emph{Mathematical Logic}, Chapter~6), with the single substitution
$\code{u}$ in place of $\ulcorner u \urcorner$.
\end{abstract}

\section{The formal system}

\subsection{Term and function syntax (Church, via Guard)}

We follow Church's basic recursive arithmetic~\cite{Church} in the
formulation of Guard~\cite{Guard1963}, with Shoenfield's
notation~\cite{Shoenfield}.  Following Shoenfield {\S}6.6, the
non-logical symbols of the theory~$T$
are a constant $\mathbf{O}$, a unary symbol $\mathbf{s}$, and finitely
many additional function symbols (Church's combinators
$\mathbf{o},\mathbf{u},\mathbf{v},\mathbf{C},\mathbf{R}$), together with
the binary identity $=$.  Variables are $\mathbf{x}_0, \mathbf{x}_1,
\ldots$ in alphabetical order, with symbol number
$\SN(\mathbf{x}_i) = 2i$.

The theory~$T$ has two further function-symbol grammars:
\begin{align*}
\mathbf{f}   &::= \mathbf{s} \mid \mathbf{o} \mid \mathbf{u} \mid \mathbf{C}(\mathbf{g}, \mathbf{f_1}, \mathbf{f_2}) \quad (\text{unary, } \mathbf{f}, \mathbf{f_1}, \mathbf{f_2} \in \Fun_1), \\
\mathbf{g}   &::= \mathbf{v} \mid \mathbf{R}(\mathbf{f}, \mathbf{g_1}, \mathbf{g_2}) \quad (\text{binary, } \mathbf{g}, \mathbf{g_1}, \mathbf{g_2} \in \Fun_2).
\end{align*}
Terms~$\mathbf{t}$ are built from $\mathbf{O}$, variables
$\mathbf{x}_i$, and applications $\mathbf{f}(\mathbf{t})$ and
$\mathbf{g}(\mathbf{t_1}, \mathbf{t_2})$.

\subsection{Formulas and derivability}

Formulas are built from equations between terms by negation and
implication:
\[
\mathbf{A}, \mathbf{B} ::= \mathbf{t_1} = \mathbf{t_2} \mid \neg \mathbf{A} \mid \mathbf{A} \supset \mathbf{B}.
\]
We write $\vdash \mathbf{A}$ for ``$\mathbf{A}$ is a theorem of~$T$''
($\mathbf{A} \in \Thmm_T$ in Shoenfield's notation).  Bare $\mathbf{A}
= \mathbf{B}$ inside $\vdash$ means $\vdash \mathbf{A} = \mathbf{B}$.

\subsection{The axioms and rules, explicitly}
\label{sec:axioms}

The Hilbert calculus is exactly Guard's Definition~7~\cite{Guard1963},
itself following Church's Princeton lectures~\cite{Church}.  The Agda
data type \verb|BRA3.Deriv.Deriv| has one constructor per Guard axiom and
rule and \emph{nothing else}; we list all fourteen axiom schemes and
three rules, in Guard's numbering, with the Agda constructor name in
brackets.  Here $\mathbf{f} \in \Fun_1$, $\mathbf{g} \in \Fun_2$,
$\mathbf{t},\mathbf{a},\mathbf{b},\mathbf{c},\mathbf{x},\mathbf{n}$ are
terms, $\mathbf{A},\mathbf{B},\mathbf{C}$ are formulas, and $a \in
\mathbb{N}$ is a variable index.

\medskip
\noindent\emph{Equational axioms (the defining equations of the
combinators).}
\begin{align*}
\textbf{0}\ [\mathtt{ax\_succ\_nonzero}]:\quad & \neg\,(\mathbf{s}(\mathbf{O}) = \mathbf{O}) \\
\textbf{1}\ [\mathtt{ax\_o}]:\quad & \mathbf{o}(\mathbf{t}) = \mathbf{O} \\
\textbf{2}\ [\mathtt{ax\_u}]:\quad & \mathbf{u}(\mathbf{t}) = \mathbf{t} \\
\textbf{3}\ [\mathtt{ax\_v}]:\quad & \mathbf{v}(\mathbf{a},\mathbf{b}) = \mathbf{b} \\
\textbf{8}\ [\mathtt{ax\_C}]:\quad & \mathbf{C}(\mathbf{g},\mathbf{f_1},\mathbf{f_2})(\mathbf{t}) = \mathbf{g}(\mathbf{f_1}(\mathbf{t}), \mathbf{f_2}(\mathbf{t})) \\
\textbf{9}\ [\mathtt{ax\_R\_base}]:\quad & \mathbf{R}(\mathbf{f},\mathbf{g_1},\mathbf{g_2})(\mathbf{x}, \mathbf{O}) = \mathbf{f}(\mathbf{x}) \\
\textbf{10}\ [\mathtt{ax\_R\_step}]:\quad & \mathbf{R}(\mathbf{f},\mathbf{g_1},\mathbf{g_2})(\mathbf{x}, \mathbf{s}\,\mathbf{n}) = \mathbf{g_1}\bigl(\mathbf{g_2}(\mathbf{x},\mathbf{n}),\ \mathbf{R}(\mathbf{f},\mathbf{g_1},\mathbf{g_2})(\mathbf{x},\mathbf{n})\bigr)
\end{align*}

\noindent\emph{Equality axioms.}
\begin{align*}
\textbf{4}\ [\mathtt{ax\_eqTrans}]:\quad & \mathbf{x} = \mathbf{y} \supset .\ \mathbf{x} = \mathbf{z} \supset \mathbf{y} = \mathbf{z} \\
\textbf{5}\ [\mathtt{ax\_eqCong_1}]:\quad & \mathbf{a} = \mathbf{b} \supset \mathbf{f}(\mathbf{a}) = \mathbf{f}(\mathbf{b}) \\
\textbf{6}\ [\mathtt{ax\_eqCong_L}]:\quad & \mathbf{a} = \mathbf{b} \supset \mathbf{g}(\mathbf{a},\mathbf{c}) = \mathbf{g}(\mathbf{b},\mathbf{c}) \\
\textbf{7}\ [\mathtt{ax\_eqCong_R}]:\quad & \mathbf{a} = \mathbf{b} \supset \mathbf{g}(\mathbf{c},\mathbf{a}) = \mathbf{g}(\mathbf{c},\mathbf{b})
\end{align*}

\noindent\emph{Propositional axioms.}
\begin{align*}
\textbf{11}\ [\mathtt{axK}]:\quad & \mathbf{A} \supset .\ \mathbf{B} \supset \mathbf{A} \\
\textbf{12}\ [\mathtt{axS}]:\quad & \bigl(\mathbf{A} \supset (\mathbf{B} \supset \mathbf{C})\bigr) \supset .\ (\mathbf{A} \supset \mathbf{B}) \supset .\ \mathbf{A} \supset \mathbf{C} \\
\textbf{13}\ [\mathtt{axNeg}]:\quad & (\neg\mathbf{A} \supset \neg\mathbf{B}) \supset .\ \mathbf{B} \supset \mathbf{A}
\end{align*}

\noindent\emph{Rules of inference.}
\begin{align*}
\textbf{I}\ [\mathtt{mp}]:\quad & \text{from } \mathbf{A} \supset \mathbf{B} \text{ and } \mathbf{A}, \text{ infer } \mathbf{B}. \\
\textbf{III}\ [\mathtt{ruleInst}]:\quad & \text{from } \mathbf{A}, \text{ infer } \substF(a, \mathbf{t}, \mathbf{A}) \quad (\text{Shoenfield's } \Sub). \\
\textbf{VI}\ [\mathtt{ruleIndNat}]:\quad & \text{from } \substF(a, \mathbf{O}, \mathbf{A}) \text{ and } \mathbf{A} \supset \substF(a, \mathbf{s}(\mathbf{x}_a), \mathbf{A}), \text{ infer } \mathbf{A}.
\end{align*}

Two economies of Guard's system, which we do \emph{not} exploit (we keep
all fourteen axioms as constructors for fidelity): Dana Scott showed
axiom~7 derivable from the others~\cite[Exercise after Def.~6]{Guard1963},
and axiom~5 is derivable from the others excluding
axiom~7~\cite[Exercise~18]{Guard1963}.  Rule~VI is stated only for
induction on $\mathbf{x}_0$ in Guard; induction on an arbitrary variable,
and the substitutivity-of-equality rule, are derivable~\cite[Exercise~19]{Guard1963}.

This list is the complete trusted base: every theorem in the
development, up to and including $\mathrm{godelII}$, is built solely from
these seventeen constructors under \texttt{-{}-safe -{}-without-K -{}-exact-split},
with no postulates.

\subsection{Numerals and the encoding $\code{u}$}

For each $n \in \mathbb{N}$ the numeral is
$\mathbf{k}_n := \mathbf{s}^n(\mathbf{O})$ (Shoenfield's
$\mathbf{k}_n$).  Guard's text uses a unary internal functor
$\num \in \Fun_1$, the \emph{numeral-to-code} map: applied to a
numeral, $\num$ produces the \emph{encoding} of that numeral.  Its
defining property is the Deriv-level closure
\[
\vdash \num(\mathbf{k}_n) = \code{\mathbf{k}_n}
\quad\text{for every meta-natural $n$,}
\]
where $\code{\mathbf{k}_n}$ is the G\"odel handle of the numeral
$\mathbf{k}_n$.  This is the BRA4 lemma $\mathit{numEq}\,n$ (its Agda
form is $\vdash \num(\mathit{natCode}\,n) = \mathit{codeTerm}(\mathit{natCode}\,n)$,
where $\mathit{natCode}\,n = \mathbf{k}_n$ and
$\mathit{codeTerm} = \code{\cdot}$ on terms).

In particular $\num$ does \emph{not} act as the identity on numerals:
for instance
\[
\num(\mathbf{s}\,\mathbf{O}) = \code{\mathbf{s}\,\mathbf{O}} = \seq{\mathbf{k}_{\mathrm{tag\_ap_1}}, \seq{\mathbf{k}_{\mathrm{tag\_s}}, \mathbf{O}}} \;\neq\; \mathbf{s}\,\mathbf{O}.
\]
The functor $\num$ is precisely Guard's underline:
$\underline{\mathbf{x}} := \num(\mathbf{x})$, the code of the value of
$\mathbf{x}$.  On a numeral $\mathbf{k}_n$ this is the encoding
$\code{\mathbf{k}_n}$; on the recursion side it satisfies the
unconditional structural closures
\[
\vdash \num(\mathbf{O}) = \mathbf{O}, \qquad
\vdash \num(\mathbf{s}\,\mathbf{t}) = \seq{\mathbf{k}_{\mathrm{tag\_ap_1}}, \seq{\mathbf{k}_{\mathrm{tag\_s}}, \num(\mathbf{t})}}
\]
(the Agda lemmas $\mathbf{num\_at\_O}$, $\mathbf{num\_at\_S}$), so that
$\num$ wraps each $\mathbf{s}$ in the same nested layout that
$\code{\cdot}$ uses for $\mathbf{s}(\cdot)$, making
$\vdash \num(\mathbf{k}_n) = \code{\mathbf{k}_n}$ a straightforward
meta-induction on~$n$.

The encoding $\code{\mathbf{u}}$ of a designator $\mathbf{u}$ is
defined by induction on the formation rules, using Shoenfield's
sequence notation $\seq{a_0, a_1, \ldots, a_n}$ ({\S}6.6).  In our
concrete Agda realisation the sequence operation is a fixed binary
functional $\Pair \in \Fun_2$ (Church's $\pi$), so that
$\seq{a, b} := \Pair(a, b)$ and longer sequences nest to the right;
the abstract notation $\seq{\cdot}$ below leaves this implementation
choice tacit.  The encoding schema is then:
\begin{align*}
\code{\mathbf{O}}              &= \mathbf{O}, \\
\code{\mathbf{x}_k}            &= \seq{\mathbf{k}_{\mathrm{tag\_var}}, \mathbf{k}_k}, \\
\code{\mathbf{f}(\mathbf{t})}  &= \seq{\mathbf{k}_{\mathrm{tag\_ap_1}}, \code{\mathbf{f}}, \code{\mathbf{t}}}, \\
\code{\mathbf{g}(\mathbf{t_1}, \mathbf{t_2})} &= \seq{\mathbf{k}_{\mathrm{tag\_ap_2}}, \code{\mathbf{g}}, \code{\mathbf{t_1}}, \code{\mathbf{t_2}}}, \\
\code{\mathbf{t_1} = \mathbf{t_2}} &= \seq{\mathbf{k}_{\mathrm{tag\_eq}}, \code{\mathbf{t_1}}, \code{\mathbf{t_2}}}, \\
\code{\neg \mathbf{A}}         &= \seq{\mathbf{k}_{\mathrm{tag\_neg}}, \code{\mathbf{A}}}, \\
\code{\mathbf{A} \supset \mathbf{B}} &= \seq{\mathbf{k}_{\mathrm{tag\_imp}}, \code{\mathbf{A}}, \code{\mathbf{B}}},
\end{align*}
with codes for the $\Fun_1$ and $\Fun_2$ symbols defined analogously.
The tags are concrete meta-naturals (e.g.\ $\mathrm{tag\_var}=1$,
$\mathrm{tag\_ap_1}=2$, $\mathrm{tag\_eq}=10$, $\mathrm{tag\_neg}=11$,
$\mathrm{tag\_imp}=12$, etc.).

\subsection{The verifier $\thmT$ and the substitution functor $\sub$}

Guard's verifier $\thmT \in \Fun_1$ (his $\mathbf{th}$, Definition~16
of~\cite{Guard1963}) is internally constructed by course-of-values
recursion so as to satisfy two structural closures:
\begin{description}
\item[Soundness-by-construction.] If $\code{\mathbf{d}}$ is the
encoding of a Hilbert derivation of~$\mathbf{A}$, then
$\vdash \thmT(\code{\mathbf{d}}) = \code{\mathbf{A}}$.

\item[Validating-decoder invariant.] On inputs that do not
encode a well-formed derivation, $\thmT$ returns the canonical truth
$\code{\mathbf{O} = \mathbf{O}}$.
\end{description}
Its construction by course-of-values recursion --- the genuinely
non-trivial step of realising such recursion inside Church's
$\{\mathbf{s},\mathbf{o},\mathbf{u},\mathbf{v},\mathbf{C},\mathbf{R}\}$
grammar --- is the subject of {\S}\ref{sec:cov}.

The substitution functor $\sub \in \Fun_2$ is Guard's Exercise~24[8]
of~\cite{Guard1963}:
\[
\vdash \sub(\mathbf{z}, \code{\mathbf{A}}) = \sbf(\seq{\mathbf{k}_0, \num(\mathbf{z})}, \code{\mathbf{A}}),
\]
where $\sbf \in \Fun_2$ is the formula-level single-variable encoded
substitution.  The point of separating $\sub$ from $\sbf$ is that
$\sub$ is the function which appears in Guard's text; $\sbf$ is the
lower-level realisation.

\subsection{The diagonal lemma}

Let $\mathbf{F}$ be the seed formula
\[
\mathbf{F} \;\equiv\; \neg\,(\thmT(\mathbf{x}_1) = \mathbf{x}_0).
\]
Define
\begin{align*}
\mathbf{H} &\;\equiv\; \Sub(\code{\mathbf{F}},\, \code{\mathbf{x}_0},\, \code{\sub(\mathbf{x}_0, \mathbf{x}_0)}), \\
i &:= \Num(\code{\mathbf{H}}), \\
\mathbf{G} &\;\equiv\; \Sub(\code{\mathbf{H}},\, \code{\mathbf{x}_0},\, \code{\mathbf{k}_i}).
\end{align*}
Then $\mathbf{G}$ is closed in $\mathbf{x}_0$ with one free variable
$\mathbf{x}_1$, and satisfies the diagonal identity
\[
\vdash \sub(\mathbf{k}_i, \mathbf{k}_i) = \code{\mathbf{G}},
\]
which is the Agda lemma \texttt{Thm14F.diag\_term\_eq}.

\section{Course-of-values recursion inside Church's recursor}
\label{sec:cov}

The verifier $\thmT$ is defined by recursion on the G\"odel number of a
derivation, where the value at a node depends on the values at
\emph{all} smaller indices (the sub-derivations).  This is
\emph{course-of-values} recursion.  Church's grammar provides only the
primitive recursor $\mathbf{R}$, and not in the textbook shape; realising
course-of-values recursion inside it was the single most delicate piece
of foundational infrastructure in the development.  Although the
technique (encode the whole history as one value) is standard in the
abstract, carrying it out in Church's specific $\mathbf{R}$ was not, for
the reasons below.

\subsection{Church's recursor is not the textbook one}

The textbook primitive recursor satisfies $\rho(\mathbf{x},\mathbf{O}) =
g(\mathbf{x})$ and $\rho(\mathbf{x},\mathbf{s}\,\mathbf{n}) =
h(\mathbf{n}, \rho(\mathbf{x},\mathbf{n}))$, feeding the recursion result
to $h$ in a fixed argument slot.  Church's $\mathbf{R}$ (axioms~9,~10 of
{\S}\ref{sec:axioms}) instead satisfies
\[
\mathbf{R}(\mathbf{f},\mathbf{g_1},\mathbf{g_2})(\mathbf{x},\mathbf{s}\,\mathbf{n}) = \mathbf{g_1}\bigl(\mathbf{g_2}(\mathbf{x},\mathbf{n}),\ \mathbf{R}(\mathbf{f},\mathbf{g_1},\mathbf{g_2})(\mathbf{x},\mathbf{n})\bigr),
\]
so the recursion result enters $\mathbf{g_1}$ as its \emph{second}
argument, while $\mathbf{g_2}(\mathbf{x},\mathbf{n})$ supplies an
auxiliary first argument.  This shape is convenient for some recursions
and awkward for others; in particular it does \emph{not} directly give
``iterate a unary $\mathbf{f}$,'' which is what course-of-values needs.

\subsection{The history-as-value encoding}

Following Church's pairing functors $J,K,L$ (Definition~6
of~\cite{Guard1963}; the realisation uses $\mathbf{Pair} = \pi$ with
projections $\mathit{Fst},\mathit{Snd}$), a finite history is encoded as a
right-nested list of pairs,
\[
\seq{v_K, v_{K-1}, \ldots, v_0} \;:=\; \mathbf{Pair}(v_K, \mathbf{Pair}(v_{K-1}, \cdots \mathbf{Pair}(v_0, \mathbf{O}))),
\]
and the recursion carries a \emph{state}
\[
\mathit{state}_K \;=\; \mathbf{Pair}(\mathbf{k}_K,\ \mathbf{Pair}(\mathit{spec},\ \mathit{table}_K)),
\qquad \mathit{table}_K = \seq{v_K, \ldots, v_0},
\]
threading the counter $\mathbf{k}_K$, an immutable \emph{spec} parameter
(recoverable as $\mathit{Fst}(\mathit{Snd}(\mathit{state}))$ at every
step, which is how the step function reads its parameter without it being
a BRA argument), and the full table of values.  At each step the new
value $v_{K+1} = \mathit{stepFun}(\mathit{Fst}\,\mathit{state},\,
\mathit{Snd}\,\mathit{state})$ is consed onto the table; the final answer
is read off the head,
\[
\thmT = \mathit{Post}\;\mathit{readOff}\;(\mathit{cov\_spec}\ \mathit{baseFun}\ \mathit{stepFun}).
\]
This is the Agda module \verb|BRA4.CoVSpec| (spec-carrying), built on
\verb|BRA3.CourseOfValues|.

\subsection{The nested-\texorpdfstring{$\mathbf{R}$}{R} trick}

The crux is that the state recursion needs to \emph{apply a unary
function to the previous state} --- i.e.\ an iteration combinator
$\mathit{iter}\ f$ with $\mathit{iter}\ f\,(\mathbf{x},\mathbf{k}) =
f^{\mathbf{k}}(\mathbf{x})$.  But Church's $\mathbf{R}$ delivers the
recursion result into $\mathbf{g_1}$'s \emph{second} slot, whereas the
naive lift $\mathbf{R}\,f\,\mathbf{v}\,\mathbf{v}$ can read it only via
the \emph{first}.  Bridging the two requires a binary $\mathbf{g_1}$ with
$\mathbf{g_1}(\mathbf{a},\mathbf{b}) = f(\mathbf{b})$ --- a function that
ignores its first argument and applies $f$ to its second.  This is itself
built as \emph{another} nested $\mathbf{R}$,
\[
\mathit{iter\_step\_fun}\,f \;=\; \mathbf{R}\,\bigl(\mathit{compose_1}\,f\,\mathbf{o}\bigr)\,\bigl(\mathbf{R}\,(\mathit{compose_1}\,f\,\mathbf{s})\,\mathbf{v}\,\mathbf{v}\bigr)\,\mathbf{v},
\]
and the defining equation $\mathbf{g_1}(\mathbf{a},\mathbf{b}) =
f(\mathbf{b})$ is then \emph{proved inside $T$ by induction}
(\texttt{ruleIndNat}, Rule~VI), not by computation.  This use of the
object-level induction rule to establish a meta-level recursion scheme is
exactly what makes the step ``non-trivial to do for Church's system'':
the iteration combinator is not primitive, it is a derived object whose
correctness is an internal theorem.

\subsection{Fuel and stability}

Because $\mathbf{R}$ recurses on a concrete numeral, $\mathit{cov\_spec}$
takes a \emph{fuel} index, and the read-off value must be shown
independent of any surplus fuel above the index actually needed --- the
\emph{stability} lemmas of \verb|BRA4.StabilityNatFuel| and
\verb|BRA4.Stability|, which establish that once the counter reaches the
target node $K$, $\mathit{readOff}(\mathit{state}_{K'}) =
\mathit{readOff}(\mathit{state}_K)$ for all $K' \ge K$.  These, together
with the per-projection preservation lemmas, are what let the
finitely-many per-tag closure equations of $\thmT$
($\thmT\_\mathrm{at}\_\mathrm{ax_0}, \ldots, \thmT\_\mathrm{at}\_\mathrm{mp},
\thmT\_\mathrm{at}\_\mathrm{sb}$, etc.) be discharged uniformly.

In sum: course-of-values recursion is \emph{representable} in BRA, but
not for free.  Its representation requires (i) a history-carrying state
threaded through Church's non-standard $\mathbf{R}$, (ii) an iteration
combinator synthesised by a nested $\mathbf{R}$ whose correctness is an
internal induction, and (iii) fuel-stability lemmas.  This is the
foundation on which $\thmT$, $\sbt$, $\sbf$, $\num$, and $\mathbf{Eval}$
(of the Berry route) all rest.

\section{The main theorem}

\begin{theorem}[G\"odel II for BRA, formalised]
\label{thm:godelII}
Let
\[
\Con_T \;\equiv\; \neg\,(\thmT(\mathbf{x}_0) = \code{\mathbf{O} = \mathbf{s}\,\mathbf{O}})
\]
be the internal consistency schema, asserting that no term codes a
derivation of $\mathbf{O} = \mathbf{s}\,\mathbf{O}$.  Then there is a
closed Hilbert-style derivation
\[
\mathrm{godelII} \;:\; \vdash \Con_T \;\Longrightarrow\; \vdash (\mathbf{O} = \mathbf{s}\,\mathbf{O}).
\]
\end{theorem}

\noindent
The Agda statement is
\begin{lstlisting}
godelII : Deriv ConSchema -> Deriv P_false
  where ConSchema = neg (eqF (ap1 thmT (var 0)) codeFalse)
        P_false   = atomic (eqn O (ap1 s O))
\end{lstlisting}
and is shipped in \verb|BRA4/Thm/Thm14GodelII.agda| (241 LoC, fresh
chain to the theorem~$7.2$\,s), with no postulates, no holes, and the
global options \verb|--safe --without-K --exact-split|.  The full BRA4
development has zero \verb|postulate| declarations anywhere in its
$\sim 50\,000$ lines.

\section{Structure of the proof}

The proof follows Guard's five-step decomposition (\emph{guard15.pdf}
p.~17), with Step~5 expressed as an explicit composition of two encoded
applications of modus ponens.  Throughout, $\mathbf{x}$ ranges over
$T$-terms at the meta level, and is \emph{universally} quantified at
the meta-level (so $\vdash \mathbf{A}(\mathbf{x})$ below stands for
``\,$\vdash \mathbf{A}$ as an open formula with free variable
$\mathbf{x}$\,'' --- giving every closed instance via Rule~III).

\paragraph{Step 1.}  By the internal Theorem~13 for unary functionals,
\[
\vdash \;\;(\thmT(\mathbf{x}) = \mathbf{k}_j) \;\supset\; (\thmT(D_\thmT(\mathbf{x})) = \code{\thmT(\num(\mathbf{x})) = \num(\mathbf{k}_j)}),
\]
where $D_\thmT \in \Fun_1$ is the Hilbert-internal handle for the
verifier $\thmT$ supplied by Theorem~12, and $j$ is the Goedel-handle
meta-natural of~$\mathbf{G}$.

\paragraph{Step 2.}  By the internal Theorem~13 for the binary
$\sub$:
\[
\vdash \;\;\thmT(D_\sub(\mathbf{k}_i, \mathbf{k}_i)) = \code{\sub(\mathbf{k}_i, \mathbf{k}_i) = \num(\mathbf{k}_j)}.
\]

\paragraph{Step 3.}  Imp-encoded transitivity of equality chains
Steps~1 and~2 through the common middle term $\num(\mathbf{k}_j)$:
\[
\vdash \;\;(\thmT(\mathbf{x}) = \mathbf{k}_j) \;\supset\; (\thmT(g(\mathbf{x})) = \code{\thmT(\num(\mathbf{x})) = \sub(\mathbf{k}_i, \mathbf{k}_i)}).
\]

\paragraph{Step 4.}  The term
\[
K_{\mathrm{part}}(\mathbf{x}) \;:=\; \seq{\mathbf{k}_{\mathrm{tag\_sb}}, \seq{\mathbf{k}_1, \num(\mathbf{x})}, \mathbf{x}}
\]
is the encoded single-variable substitution that, under the diagonal
identity $\sub(\mathbf{k}_i, \mathbf{k}_i) = \code{\mathbf{G}}$,
satisfies
\[
\vdash \;\;(\thmT(\mathbf{x}) = \mathbf{k}_j) \;\supset\; (\thmT(K_{\mathrm{part}}(\mathbf{x})) = \code{\neg\,(\thmT(\num(\mathbf{x})) = \sub(\mathbf{k}_i, \mathbf{k}_i))}).
\]

\paragraph{Step 5a.}  Let $t' := \code{\mathbf{axExFalso}\,\mathbf{A}\,\bot}$
be the encoded derivation of Guard's schema $\mathbf{A} \supset
(\neg \mathbf{A} \supset \bot)$ with $\mathbf{A} \equiv (\mathbf{x}_0
= \mathbf{x}_1)$ and $\bot \equiv (\mathbf{O} = \mathbf{s}\,\mathbf{O})$.
Set $h(\mathbf{x})$ to be \emph{two nested single-variable substitution
wraps} of $t'$ (var~$0$ outermost, var~$1$ innermost):
\[
h(\mathbf{x}) \;:=\; \seq{\mathbf{k}_{\mathrm{tag\_sb}},\, \seq{\mathbf{k}_0, \code{\thmT(\num(\mathbf{x}))}},\,
   \seq{\mathbf{k}_{\mathrm{tag\_sb}},\, \seq{\mathbf{k}_1, \code{\sub(\mathbf{k}_i, \mathbf{k}_i)}},\, t'}}.
\]
$\thmT$ decodes the outer wrap by its single $\mathbf{sb}$-clause into
$\sbf(\seq{\mathbf{k}_0, \code{\thmT(\num\,\mathbf{x})}},\, \cdot)$ applied to
the inner wrap, which decodes likewise into
$\sbf(\seq{\mathbf{k}_1, \code{\sub(\mathbf{k}_i,\mathbf{k}_i)}},\, \thmT(t'))$.
Both substituents are $\num$-based (var-free) codes, hence inert under the
outer re-scan ({\S}\ref{sec:sb2}); writing
$E := \thmT(\num(\mathbf{x})) = \sub(\mathbf{k}_i, \mathbf{k}_i)$,
\[
\vdash \;\;\thmT(h(\mathbf{x})) = \code{E \;\supset\; \neg E \;\supset\; (\mathbf{O} = \mathbf{s}\,\mathbf{O})}.
\]

\paragraph{Step 5b.}  Two applications of the Carneiro-lifted encoded
modus ponens $\mathbf{imp\_encoded\_mp}$ compose Steps~3, 4, and~5a to
produce
\[
\vdash \;\;(\thmT(\mathbf{x}) = \mathbf{k}_j) \;\supset\; (\thmT(\mathit{bigterm}(\mathbf{x})) = \code{\mathbf{O} = \mathbf{s}\,\mathbf{O}}),
\]
where
\[
\mathit{bigterm}(\mathbf{x}) \;:=\; \seq{\mathbf{k}_{\mathrm{tag\_mp}}, \seq{\mathbf{k}_{\mathrm{tag\_mp}}, h(\mathbf{x}), g(\mathbf{x})}, K_{\mathrm{part}}(\mathbf{x})}.
\]

\paragraph{Final assembly.}  Instantiate the universal Step~5 at
$\mathbf{x} = \mathbf{x}_1$.  From $\vdash \Con_T$ by Rule~III applied
to $\mathit{bigterm}(\mathbf{x}_1)$ we obtain
\[
\vdash \;\neg\,(\thmT(\mathit{bigterm}(\mathbf{x}_1)) = \code{\mathbf{O} = \mathbf{s}\,\mathbf{O}}).
\]
The classical contrapositive $\mathbf{axContrapos}$ then yields
\[
\vdash \;\neg\,(\thmT(\mathbf{x}_1) = \mathbf{k}_j).
\]
The predicate-Leibniz rule $\mathbf{substF\_cong}$ applied at the
formula
\[
\Phi \;\equiv\; \neg\,(\thmT(\mathbf{x}_1) = \mathbf{x}_2)
\]
with the equation $\vdash \sub(\mathbf{k}_i, \mathbf{k}_i) = \mathbf{k}_j$
converts this to $\vdash \mathbf{G}$.  An inlined re-derivation of
Theorem~11 (G\"odel~I) for the sub-form of $\mathbf{G}$ then yields
$\vdash (\mathbf{O} = \mathbf{s}\,\mathbf{O})$.

\section{Implicit facts in Guard's text, made explicit}

The formalisation forced a number of small but non-trivial corrections,
disambiguations, or additions to Guard 1963.

\subsection{Underline notation: the internal $\num$ versus the meta-encoding $\code{\cdot}$}

Guard's text uses an underline notation $\underline{\mathbf{x}}$ whose
meaning shifts between (i)~the internal $T$-term $\num(\mathbf{x})$
and (ii)~the meta-level G\"odel handle $\code{\mathbf{x}}$.  These are
genuinely different objects:
\begin{itemize}
\item $\num(\mathbf{x})$ is a $T$-term --- the unary functor $\num \in
\Fun_1$ applied to $\mathbf{x}$.  It \emph{computes}, inside $T$, the
encoding of the value of $\mathbf{x}$.
\item $\code{\mathbf{x}}$ is the meta-level encoding of the
\emph{syntactic} object $\mathbf{x}$ (the G\"odel handle), a closed
term computed outside $T$.
\end{itemize}
They are connected only by the Deriv-level closure on closed numerals:
\[
\vdash \num(\mathbf{k}_n) = \code{\mathbf{k}_n}
\quad\text{(the lemma $\mathit{numEq}\,n$),}
\]
and they differ for non-numeral $\mathbf{x}$ (e.g.\ for a variable
$\mathbf{x}_0$, the term $\num(\mathbf{x}_0)$ is open whereas
$\code{\mathbf{x}_0} = \seq{\mathbf{k}_{\mathrm{tag\_var}}, \mathbf{k}_0}$
is a fixed closed handle).  Note that \emph{neither} of these equals
$\mathbf{x}$ itself: in particular $\num(\mathbf{s}\,\mathbf{O}) =
\code{\mathbf{s}\,\mathbf{O}} \neq \mathbf{s}\,\mathbf{O}$, so the
underline is never the identity.

In the Agda development the file \verb|Thm14F.agda| introduces the
\emph{slot abbreviation} $\mathrm{code}\,\mathbf{t} := \num(\mathbf{t})$
(its $\mathrm{code}$ is literally $\num$), used in the encoded
positions where Guard writes $\underline{\mathbf{x}}$.  This
abbreviation coincides with the meta-encoding $\code{\cdot}$ only on
closed numerals --- precisely via $\mathit{numEq}$ --- and the
formalisation makes every such bridge explicit
(e.g.\ $\vdash \num(\mathbf{k}_i) = \code{\mathbf{k}_i}$ at the
diagonal numeral $i$).

\subsection{The two G's of the diagonal lemma}

There are two distinct formulas, both called ``$\mathbf{G}$'' in
informal treatments of the diagonal lemma:
\begin{align*}
\mathbf{G}_{\sub}  &:\equiv\; \neg\,(\thmT(\mathbf{x}_1) = \sub(\mathbf{k}_i, \mathbf{k}_i)), \\
\mathbf{G}_{\sbf}  &:\equiv\; \neg\,(\thmT(\mathbf{x}_1) = \sbf(\seq{\mathbf{k}_0, \num(\mathbf{k}_i)}, \mathbf{k}_i)).
\end{align*}
These are equal in~$T$ only up to the closure $\mathbf{sub\_eq}$,
which is a term-level $T$-derivable equation, \emph{not} a
Formula-syntactic equality.  Consequently, a proof of G\"odel~I phrased
for $\mathbf{G}_{\sbf}$ (as in our \verb|BRA4.GoedelI|) does not
literally apply to $\mathbf{G}_{\sub}$.  The formalisation bridges this
gap either by inlining a copy of G\"odel~I for $\mathbf{G}_{\sub}$ (the
route taken in \verb|Thm14GodelII|), or by invoking the
\emph{predicate-Leibniz} rule $\mathbf{substF\_cong}$ to transfer
derivability across the term-equation $\sub(\mathbf{k}_i, \mathbf{k}_i)
= \mathbf{k}_j$.  Guard's text treats the two $\mathbf{G}$'s as
interchangeable without comment.

\subsection{Modus-ponens slot order}

{\sloppy Guard writes \emph{mp}(antecedent, implication) in
his encoded-derivation tables, while our
$\seq{\mathbf{k}_{\mathrm{tag\_mp}}, \mathrm{impl}, \mathrm{ant}}$ uses
the opposite order.\par}  Guard's expression for
$\mathit{bigterm}(\mathbf{x})$,
\[
4J\bigl[\,4J(J(\num\,\mathbf{x},1),\, \mathbf{x})+1,\ \ 4J(g(\mathbf{x}),\, h(\mathbf{x}))+2\,\bigr]+2,
\]
thus hides which sub-slot is which, and a careful sign-by-sign
translation is needed.

\subsection{Closedness side conditions for Rule III}

Rule III ($\substT, \substF$) only reduces $\substT(\mathbf{x}_a,
\mathbf{t}, \mathbf{t}')$ when $\mathbf{x}_a$ and $\mathbf{t}'$ match
by definitional pattern.  For closed sub-terms (numerals, the encoded
false formula $\code{\mathbf{O} = \mathbf{s}\,\mathbf{O}}$, the
diagonal numeral $\mathbf{k}_i$) the formalisation requires explicit
$\Closed$ witnesses and bridges via \texttt{closedAt}.  Guard's text
uses ``\,$\mathbf{k}_i$ is a numeral so its substitution is
trivial\,''.  The Agda development must construct, e.g.\
\[
\mathit{closed\_codeFalse} \;:\; \Closed\,\code{\mathbf{O} = \mathbf{s}\,\mathbf{O}}
\]
as a nested $\seq{\cdot}$/numeral proof tree.

\subsection{The validating-decoder invariant on $\thmT$}

Guard's verifier $\thmT$ is described informally as ``returns the code
of the proved formula''.  But for inputs that are \emph{not} the
encoding of a derivation, $\thmT$ must still be defined and must still
satisfy the closure equations used in the proof (e.g.\
$\vdash \thmT(\code{\mathrm{mp}(p,q)}) = \cdots$).  We adopt the
convention
\[
\vdash \thmT(\mathit{garbage}) = \code{\mathbf{O} = \mathbf{O}},
\]
(the canonical trivial truth), which makes the closure equations
$\mathbf{thmT\_at\_mp}$, $\mathbf{thmT\_at\_sb}$, etc.\ universal in
their raw $\Pair$-encoded inputs.  This invariant is implicit in Guard
but essential to make the closures usable.

\subsection{The encoded $t'$ schema}

Guard writes ``\,$t'$ is the encoded derivation of $\mathbf{x}_0 =
\mathbf{x}_1 \supset \mathbf{x}_0 \neq \mathbf{x}_1 \supset \mathbf{O}
= \mathbf{s}\,\mathbf{O}$\,''.  The formalisation exhibits this
concretely as
\[
t' \;=\; \code{\mathbf{axExFalso}\,(\mathbf{x}_0 = \mathbf{x}_1)\,(\mathbf{O} = \mathbf{s}\,\mathbf{O})}
\]
with $\mathbf{axExFalso}$ built from $\mathbf{axK}, \mathbf{axS},
\mathbf{axNeg}$, and a completeness handle
$\mathbf{thmT\_complete\_rec}$ proving the verifier's expected output.
Guard glosses both pieces.

\subsection{Carneiro lift for hypothetical derivations}
\label{sec:carneiro}

Guard frequently asserts implications of the form ``\,under the
hypothesis $\mathbf{P}$ we have $\mathbf{A}$ and $\mathbf{B}$, hence
$\mathbf{C}$\,''.  In a pure Hilbert calculus there is no
implication-introduction rule, and our \verb|Deriv| is
\emph{context-free} ($\verb|Deriv| : \mathit{Formula} \to \mathit{Set}$,
no hypothesis context), so the only way to express ``under hypothesis
$\mathbf{P}$'' is the object formula $\mathbf{P} \supset (\cdots)$, and
the only way to \emph{compose} such conditionals is to lift every
ordinary combinator to its $\mathbf{P} \supset (\cdots)$ form via
Carneiro's $\mathbf{axS}/\mathbf{axK}$ technique~\cite{Carneiro}.  The
formalisation extracts this into a small toolkit in
\verb|BRA4.Thm12.ImpHelpers|:
\[
\mathbf{impLift},\ \mathbf{impMp},\ \mathbf{impCong_1},\ \mathbf{impCong_L},\ \mathbf{impCong_R},\ \mathbf{impRuleSym},\ \mathbf{impEqTrans},
\]
together with the encoded modus ponens $\mathbf{imp\_encoded\_mp}$
(\verb|BRA4.Thm12.EncodedMp|).

\paragraph{Is it essential?}  Relative to the context-free
\verb|Deriv| representation, yes: the theorems of the Theorem~14
argument are intrinsically conditional --- Step~5 is literally
$\mathrm{th}(\mathbf{x}) = j \supset \mathrm{th}(\mathit{bigterm}(\mathbf{x}))
= \code{\mathbf{O}=\mathbf{s}\,\mathbf{O}}$ --- and this conditional is
exactly what is combined with $\Con_T$ to produce $\mathbf{G}$, so it
cannot be restructured away.  It is \emph{not} essential in the absolute
sense: a deduction theorem for BRA is admissible (\cite[Exercise~19]{Guard1963}),
so one could instead carry a hypothesis context in \verb|Deriv| and prove
the deduction theorem once.  The Carneiro lift is the \emph{pointwise}
deduction theorem --- lift the combinators actually used, rather than
prove the general metatheorem --- and given our context-free judgement it
is the lightweight, effectively forced choice.

\paragraph{A concrete example where it helped.}  The step in
$\thmT$'s Step~4 (\verb|BRA4.Thm.Thm14Step4|) that \emph{uses the
G\"odel hypothesis itself}.  Under $\mathbf{P} = (\mathrm{th}(\mathbf{x})
= j)$ one must rewrite $\sbf(\mathit{cSpec}, \mathrm{th}(\mathbf{x}))$
into $\sbf(\mathit{cSpec}, j)$ --- i.e.\ apply the assumption
``$\mathrm{th}(\mathbf{x}) = j$'' \emph{inside} a congruence, without
discharging it:
\begin{lstlisting}
hyp_imp : Deriv (imp P (eqF (ap1 thmT x) j))
hyp_imp = impRefl P             -- assuming P, P holds: i.e. assuming th(x)=j, th(x)=j

step_B_imp : Deriv (imp P (eqF (ap2 sbf cSpec (ap1 thmT x)) (ap2 sbf cSpec j)))
step_B_imp = impCongR {P} sbf (ap1 thmT x) j cSpec hyp_imp
\end{lstlisting}
Here $\mathbf{impRefl}\,\mathbf{P}$ extracts the hypothesis as
$\mathbf{P} \supset \mathbf{P}$ --- and $\mathbf{P}$ \emph{is} the
equation $\mathrm{th}(\mathbf{x}) = j$ --- while $\mathbf{impCong_R}$
threads it through the congruence of $\sbf$, yielding
$\mathbf{P} \supset (\sbf(\mathit{cSpec}, \mathrm{th}(\mathbf{x})) =
\sbf(\mathit{cSpec}, j))$.  Unfolded, this is just $\mathbf{ax\_eqCong_R}$
wrapped with $\mathbf{axS}/\mathbf{axK}$; the toolkit makes it one line.

The same proof shows the lift is doing genuine work of two distinct
kinds at once.  The surrounding steps A, C, D of Step~4 are
\emph{unconditional} facts (the closure $\thmT\_\mathrm{at}\_\mathrm{sb}$,
the diagonal bridge $\mathit{codeFormulaG\_eq\_j}$, the per-shape
$\sbf/\sbt$ evaluations); each is carried past $\mathbf{P}$ by
$\mathbf{impLift}$ (a bare $\mathbf{axK}$ wrap).  Only step~B
\emph{consumes} the hypothesis ($\mathbf{impRefl}\,\mathbf{P}$ +
$\mathbf{impCong_R}$).  And $\mathbf{impEqTrans}$ chains all four equalities
under the single $\mathbf{P}$ to conclude
$\mathbf{P} \supset (\mathrm{th}(K_{\mathrm{part}}(\mathbf{x})) =
\code{\neg(\mathrm{th}(\underline{\mathbf{x}}) = \sub(\mathbf{k}_i,\mathbf{k}_i))})$.
Carrying unconditional facts past $\mathbf{P}$ and applying $\mathbf{P}$
where needed is exactly the bookkeeping a deduction theorem would
otherwise perform; the Carneiro lift does it pointwise.  Its headline
consumer is $\mathbf{imp\_encoded\_mp}$: the two encoded modus-ponens
steps of Step~5b are composed entirely under $P_x$ via
$\mathbf{imp\_thmT\_at\_mp} + \mathbf{impEqTrans}$.  Guard never names the
principle.

\subsection{Predicate-Leibniz substitutivity}

Guard's transition from
``\,$\neg\,(\thmT(\mathbf{x}_1) = \mathbf{k}_j)$\,''
to
``\,$\neg\,(\thmT(\mathbf{x}_1) = \sub(\mathbf{k}_i, \mathbf{k}_i))$\,''
is stated as an obvious rewrite.  In a Hilbert calculus this requires
the \emph{predicate-Leibniz} principle
\[
\vdash \mathbf{a} = \mathbf{b} \;\Longrightarrow\; \vdash \Sub(\code{\Phi}, \code{\mathbf{x}_k}, \code{\mathbf{a}}) \supset \Sub(\code{\Phi}, \code{\mathbf{x}_k}, \code{\mathbf{b}}),
\]
which is itself a small but non-trivial Hilbert-meta-theorem (proved by
induction on $\Phi$, using $\mathbf{axContrapos}$ on the negative
slot; see \verb|BRA3.Substitutivity|).  Guard does not mention it.

\subsection{Three small typographical issues}

\begin{itemize}
\item Guard writes ``$j$'' in places where his own underline convention
requires $\underline{j}$ (= $\num\,\mathbf{k}_j = \code{\mathbf{k}_j}$).
The formalisation makes the distinction explicit at every occurrence.

\item Definition~12 conflates the underline notation with the
numeral.  The formalisation requires the precise statement that
$\vdash \num(\mathbf{k}_n) = \code{\mathbf{k}_n}$ \emph{for every
meta-natural~$n$} --- i.e.\ $\num$ maps a numeral to its
\emph{encoding}, not to itself --- instantiated through the lemma
$\mathit{numEq}$.  In particular $\num(\mathbf{s}\,\mathbf{O}) =
\code{\mathbf{s}\,\mathbf{O}} \neq \mathbf{s}\,\mathbf{O}$.  The text
leaves this implicit and the underline notation makes it easy to
mistake $\num(\mathbf{x})$ for $\mathbf{x}$ itself.

\item In the description of $h(\mathbf{x})$, Guard does not commit to a
representation of the substitution.  The formalisation realises the
two-variable effect as two \emph{nested single-variable} wraps
$\seq{\mathbf{k}_{\mathrm{tag\_sb}}, \seq{\mathbf{k}_0, S_0},
\seq{\mathbf{k}_{\mathrm{tag\_sb}}, \seq{\mathbf{k}_1, S_1}, t}}$ and uses
the numeral-inertness lemma ({\S}\ref{sec:sb2}) to pass the inner
substituent $S_1$ through the outer re-scan.
\end{itemize}

\section{Single-variable substitution suffices: numeral-inertness}
\label{sec:sb2}

The verifier's substitution clause is the \emph{single-variable} encoded
substitution $\sbf/\sbt$ (cf.\ Guard's $\mathbf{sb}$, Def.~16).  Both the
compound-functor congruences feeding Theorem~12 and the two-variable wrap of
Step~5a require, on the face of it, a \emph{simultaneous} substitution of
several variables at once.  We obtain all of them from \emph{nested
single-$\mathbf{sb}$ wraps}, with no dedicated multi-variable functor.  The
one fact that makes this work is a numeral-inertness lemma; we record it
here because it is exactly the place Guard's text leaves the argument
implicit, and because the alternative (a genuine simultaneous-substitution
primitive) carries a different proof-theoretic commitment relevant to one of
Guard's open problems.

\paragraph{The apparent obstruction.}  A nested wrap substitutes the
variables one at a time, so an outer pass re-scans a substituent already
plopped by an inner pass.  In Theorem~12's compound cases the substituents
are $\num\,X$ for the \emph{object} variable $X$, and for open $X$ (e.g.\
$X=\mathbf{x}_1$) the term $\num(\mathbf{x}_1)$ does not reduce:
$\num = \mathbf{C}\,\mathit{numAux}\,\mathbf{o}\,\mathbf{u}$ with
$\mathit{numAux}=\mathbf{R}\,\mathbf{o}\,\mathit{numStep}\,\mathbf{v}$, whose
$\mathbf{R}$-recursion fires only when its second argument is syntactically
$\mathbf{O}$ or $\mathbf{s}(\cdot)$; a variable is neither, so
$\num(\mathbf{x}_1)$ stays a stuck open application.  Since $\sbt$ dispatches
on $\mathit{get\_tag}(\mathit{input})$ and that tag is an \emph{open}
expression on the stuck term, naive evaluation of
$\sbt(\seq{\mathbf{k}_k, S}, \num(\mathbf{x}_1))$ never settles.  Read
\emph{computationally}, the re-scan stalls --- which is why a simultaneous
plop (no re-scan) looks unavoidable.

\paragraph{The numeral-inertness lemma.}  The obstruction conflates
computation with \emph{derivability}.  We never need $\sbt$ to \emph{compute}
on a stuck $\num(\mathbf{x}_1)$; we need the \emph{equation}
\[
\vdash\ \sbt(\seq{\mathbf{k}_k, S},\, \num\,X) \;=\; \num\,X
\qquad\text{for an \emph{arbitrary} object term } X
\]
(the Agda lemma $\mathbf{sbt\_num\_inert}$, \verb|BRA4.NumInert|), and this is
a \emph{theorem} of the object system, proved by \emph{internal} induction on
$X$ (Rule~III / $\mathbf{ruleIndNat}$), not by inspection of syntax.  The
base case is $\sbt(\cdots, \num\,\mathbf{O}) = \sbt(\cdots,\mathbf{O}) =
\mathbf{O} = \num\,\mathbf{O}$; the step uses $\num(\mathbf{s}\,X) =
\seq{\mathbf{k}_{\mathrm{tag\_ap_1}}, \seq{\mathbf{k}_{\mathrm{tag\_s}},
\num\,X}}$ (the lemma $\mathbf{num\_at\_S}$), which at \emph{each} inductive
step rewrites the numeral to a $\Pair$-code the descent can enter,
discharging the recursive call by the induction hypothesis.  So the descent
never meets a bare $\num(\mathbf{x}_1)$: it meets a planted num-leaf and
passes through it.  Lifting through the $\mathbf{ap_1}/\mathbf{ap_2}$
code-wrappers (the predicate $\mathit{NumCode}$ in \verb|BRA4.SbStep|),
\emph{every} substituent built from numeral leaves is inert under
\emph{every} single-variable substitution.

\paragraph{Intuition: the internalisation of ``$\mathbf{S}^n\mathbf{0}$ is
closed''.}  The lemma says nothing more than the elementary fact that a
numeral is a \emph{closed} term: $\mathbf{S}^n\mathbf{0}$ contains no
variable, so every substitution acts as the identity on it.  For a concrete
numeral this is immediate --- its code is a var-free $\code{\cdot}$-tree with
no $\seq{\mathbf{k}_{\mathrm{tag\_var}}, \cdot}$ leaf, so the encoded re-scan
visibly passes through.  The only subtlety, and the reason an internal
induction is needed at all, is the \emph{generic} statement: for a variable
$X$ the term $\num\,X$ is a \emph{stuck open} application, not literally a
numeral, so its closedness is invisible to the tag-dispatch and must be
\emph{promoted to an internal theorem}.  $\mathbf{ruleIndNat}$ is precisely
that promotion: it internalises ``$\num\,X$ is (the code of) a closed
numeral, for all $X$'' as a derivable equation.

\paragraph{Consequence: the multi-variable machinery is eliminated.}  With
inertness in hand, each compound-functor congruence of Theorem~12 --- e.g.\
specialising the binary congruence at the \emph{open} terms
$g(\underline{\mathbf{x}}), g(\mathbf{x}), h(\underline{\mathbf{x}}),
h(\mathbf{x})$ in building $D_{\mathbf{C}(f,g,h)}$ --- and the Step~5a wrap of
Theorem~14 are realised by \emph{nested single-$\mathbf{sb}$ wraps}: an inner
pass plants a substituent, an outer pass plants the next and re-scans the
inner one, which inertness leaves fixed.  Concretely $\thmT$ decodes a
$k$-fold nested $\seq{\mathbf{k}_{\mathrm{tag\_sb}}, \cdots}$ wrap by $k$
applications of its single $\mathbf{sb}$-clause into
$\sbf(\mathrm{spec}_0,\sbf(\mathrm{spec}_1,\cdots))$.  The dedicated two- and
three-variable functors $\sbttwo,\sbftwo,\mathbf{sbt_3},\mathbf{sbf_3}$, their
contract records, and the two extra discriminator levels in the $\thmT$
cascade --- roughly ten thousand lines across a dozen files --- are thereby
\emph{removed}.  This recovers exactly Guard's own explicit Theorem~14
construction (guard15.pdf p.~17), where the multi-variable effect is obtained
by \emph{nested single-$\mathbf{sb}$ steps with the open substituent applied
last}; what Guard left implicit --- that the re-scanned substituents
genuinely pass through --- is the inertness lemma.

\paragraph{The induction-avoiding alternative, and an open problem.}  There
is a second, equally sound route that we deliberately keep on record.  A
genuine \emph{simultaneous} substitution functor $\subtwo/\mathbf{sub_3}$
plops all substituents at once; nothing is re-scanned, so its closure proofs
need \emph{no} internal induction.  (It is sound: its encoded closures are
the encoded-level counterpart of the \emph{derived admissible} rule
$\mathbf{ruleInst2}$, $k_1\neq k_2 \Rightarrow {\vdash}\mathbf{P} \Rightarrow
{\vdash}\mathit{simSubstF}(k_1,t_1,k_2,t_2,\mathbf{P})$ --- Church's Standard
Metatheorem~VII, \emph{fully derived} from single substitution by
$\alpha$-renaming through a fresh variable in \verb|BRA3.RuleInst2| --- so it
adds no proof-theoretic strength.)  The contrast is precise: the
\emph{abstract} rule reduces to single substitution on raw formulas, but the
\emph{encoded} simultaneous functor does not reduce to iterated encoded
$\sbt$ \emph{by computation} (the re-scan stalls) --- only \emph{by
derivability}, and that derivation is exactly $\mathbf{sbt\_num\_inert}$,
which uses internal induction.  So the two routes differ in \emph{where they
spend induction}: the inertness route puts an internal induction into the
verifier's substitution-correctness; the simultaneous-substitution route
keeps that piece induction-free.

This bookkeeping is not idle.  Guard (guard19.pdf, p.~19, Def.~14 and the
\textsc{open problem} following Theorem~18) defines the \emph{induction-free}
fragment $\mathrm{BRA}_0$ --- the variable-free instances of BRA's axioms
with \emph{modus ponens as the only rule} --- whose verifier $\mathbf{th}_0$
has \emph{only} an axiom clause and an mp clause (no substitution clause, no
induction clause), and \emph{conjectures} that its consistency
$\mathbf{th}_0(y) \neq \code{\mathbf{O}=\mathbf{s}\,\mathbf{O}}$, though
valid, is \emph{unprovable in BRA} --- a putative counterexample to the usual
provable consistency of induction-free fragments.  Whether, and where, a
verifier's substitution machinery can be made induction-free is exactly the
kind of distinction that question turns on; the inert-lemma versus
simultaneous-substitution choice above is one concrete instance, which is why
we keep the simultaneous functors on record even though the shipped
development uses the single-variable route.

\section{Conclusion}

The formalisation closes Guard's argument for G\"odel's second
incompleteness theorem for BRA in a completely concrete,
constructively-verified form, and clarifies the implicit mathematical
machinery (closedness side conditions, predicate-Leibniz, Carneiro lift,
validating-decoder invariant, the two-forms-of-$\mathbf{G}$ distinction) on
which Guard's text relies without comment.  As {\S}\ref{sec:sb2} documents,
the multi-variable effects needed in Theorems~12 and~14 are obtained from
\emph{nested single-variable substitution}, the re-scan of each
numeral-based substituent being discharged by the numeral-inertness lemma
$\mathbf{sbt\_num\_inert}$ --- the internalisation, by Rule~III, of the fact
that a numeral is closed.  This recovers Guard's own explicit construction
(guard15.pdf p.~17) and eliminates the dedicated two- and three-variable
simultaneous-substitution machinery.  We keep on record, as a remark bearing
on Guard's open problem about the consistency of the induction-free fragment
$\mathrm{BRA}_0$ (guard19.pdf p.~19), that the simultaneous-substitution
route is the \emph{induction-free} alternative to the inertness lemma: the
two agree on the encoded result but differ on whether the verifier's
substitution-correctness is proved with or without internal induction.

\paragraph{Acknowledgements.}  {\sloppy The Carneiro-lift toolkit
follows Mario Carneiro's natural-deduction note~\cite{Carneiro}.  The
course-of-values infrastructure (the Agda modules
\texttt{BRA3.CourseOfValues}, \texttt{BRA3.PairAlgebra}, and
\texttt{BRA3.RuleInst2}) builds on Church's system as transmitted by
Guard~\cite{Guard1963,Church}.\par}

\begin{thebibliography}{9}

\bibitem{Church}
A.~Church.
\emph{Basic recursive arithmetic} (the system of functors
$\mathbf{s},\mathbf{o},\mathbf{u},\mathbf{v},\mathbf{C},\mathbf{R}$ and
the axioms of {\S}\ref{sec:axioms}).
Lectures, Logic Seminar, Princeton University, 1957--1959.
The G\"odel numbering and the first incompleteness theorem for the
system are, in Guard's words, ``essentially those of
Church''~\cite[p.~13]{Guard1963}.

\bibitem{Guard1963}
R.~B.~Guard.
\emph{Lecture Notes on Recursive Arithmetic}.
1963.
(The working excerpt used here, \texttt{guard15.pdf}, contains
Definitions~6--16, Exercises~17--24, and Theorems~8--16, including the
verifier $\mathbf{th}$ (Def.~16) and G\"odel's first and second theorems
for BRA (Thms.~11 and~14).)

\bibitem{Shoenfield}
J.~R.~Shoenfield.
\emph{Mathematical Logic}.
Addison--Wesley, 1967.
(Chapter~6, {\S}{\S}6.6--6.7: expression numbers and representability;
the source of the $\code{u}$/$\seq{\cdot}$/$\Sub$/$\Num$/$\mathbf{th}$
notation followed throughout.)

\bibitem{Carneiro}
M.~Carneiro.
\emph{Natural Deduction in the Metamath Proof Language}.
\url{https://us.metamath.org/ocat/natded.pdf}.
The $\mathrm{axS}/\mathrm{axK}$ technique behind
\texttt{BRA4.Thm12.ImpHelpers} and \texttt{imp\_encoded\_mp}.

\end{thebibliography}

\end{document}
